\subsection{Motivation}
\label{subsec:motivation}
Urban charging demand is highly uneven across time and location, which might cause queues at highly frequented stations while other sites remain virtually unused. At the same time, cars that are parked for extended periods of time (in residential or commercial areas) can remain plugged in, creating opportunities for bidirectional energy exchange. The spatial distribution of charging stations therefore affects both user experience and grid load. Understanding the spatial and temporal patterns of charging demand is a prerequisite for infrastructure planning. However, there is a lack of open, reproducible methods for assessing candidate locations based on real-world mobility patterns and grid constraints.
 